\paragraph{C.3i (Refuted in every proposed degree-one range): strict local Picard restriction without a depth hypothesis.}
For every $N\ge4$ there is an excellent strictly Henselian
normal local complex ring $R$ of dimension $N$, and a
nonzero $f\in\mathfrak m_R$, such that
\[
 X=\operatorname{Spec}R\setminus\{\mathfrak m_R\},\qquad
 Y=X\cap V(f)
\]
are regular, while
\begin{equation}\label{eq:C3-picard-kernel}
 \ker\bigl(\operatorname{Pic}(X)\longrightarrow
                  \operatorname{Pic}(Y)\bigr)
             \simeq (\mathbb C,+)^3.
\end{equation}
The map is surjective in this example. Thus the assertion
on 14:107 that $A^1(X)\to A^1(Y)$ is injective for $N\ge4$,
and is an isomorphism for $N\ge5$, is false on the
rational-equivalence reading $A^1=\mathrm{CH}^1=\operatorname{Pic}$.
The failure persists after tensoring with $\mathbb Q$ and
after discarding finite groups.

\begin{proof}
Put $m=N-2\ge2$ and take the fixed smooth plane quartic
\[
 C:\ z_0^4+z_1^4+z_2^4=0\subset\mathbb P^2_{\mathbb C}.
\]
Its partial derivatives have no simultaneous projective
zero, and adjunction gives $K_C=\mathcal O_C(1)$ and
genus $3$. Set
\[
 W=\mathbb P^m\times C,\qquad
 L=\mathcal O_{\mathbb P^m}(1)\boxtimes K_C,
 \qquad f=x_0z_0+x_1z_1+x_2z_2\in H^0(W,L).
\]
Here unused coordinates $x_3,\ldots,x_m$ are allowed.
The morphism
\[
 \mathcal O_C^{m+1}\longrightarrow K_C,
       \qquad (a_0,\ldots,a_m)\longmapsto
                         a_0z_0+a_1z_1+a_2z_2
\]
is surjective. Its kernel $E$ is a vector bundle of rank
$m$, and $H=V_W(f)$ is its projective bundle of lines.
Consequently $H$ is smooth and connected, of dimension
$m$. Both the source and target varieties used below
therefore have dimension at least two.

The full section ring
\[
 B=\bigoplus_{a\ge0}H^0(W,L^a)
\]
is a normal domain, generated in degree one. Indeed it
is the Segre product of a polynomial ring with the
homogeneous coordinate ring of the plane quartic.
The complete section rings of these projectively normal
embeddings give the asserted normality and generation.
Let $v$ denote the cone vertex, let
$A=B_{B_+}$, and set $R=A^{\mathrm{sh}}$ with the
residue field $\mathbb C$ fixed. This agrees with the
henselization, since the residue field is algebraically
closed. The ring is excellent, normal and of dimension
$m+2=N$. The punctured affine cone over $(W,L)$ is the
complement of the zero section in $\operatorname{Tot}(L^{-1})$.
It is smooth, and the section $f$ cuts out the corresponding
punctured line bundle over $H$. Passing to the strict
henselization preserves regularity, so $X,Y$ have the
required regularity.

We first check that the possible nonnormality of the
closed hyperplane cone at its vertex is harmless here.
Let
\[
 B_H=\bigoplus_{a\ge0}H^0(H,L_H^a),\qquad L_H=L|_H.
\]
The exact sequence
\[
 0\longrightarrow L^{a-1}\xrightarrow{\ f\ }L^a
        \longrightarrow L_H^a\longrightarrow0
\]
gives an injection $B/(f)\hookrightarrow B_H$ whose
cokernel is supported in finitely many degrees: the
next coherent cohomology vanishes for $a$ sufficiently
large. Thus $B_H$ is a finite $B/(f)$-module. It is a
normal domain and the map is an isomorphism away from
the vertex. In particular it is the normalization and
the $S_2$-hull of the hyperplane cone; it has just one
point above the vertex, with residue field $\mathbb C$.
Finite base change and henselization commute in this
one-point situation. Moreover henselization commutes
with quotient by $f$. It follows that the puncture
$Y$ is also the puncture of the strict henselization of
$(B_H)_{(B_H)_+}$. We may compute its local Picard group
with this normal cone. No depth hypothesis on $R/(f)$
has been inserted.

Here is the actual group calculation that will be used.
In characteristic zero, for a smooth projective variety
$V$ of dimension at least two and an ample embedding
line bundle $M$, the local Picard \emph{scheme} of its
cone has the natural exact sequence
\begin{equation}\label{eq:C3-Kollar-cone}
 0\longrightarrow\bigoplus_{a>0}H^1(V,M^a)
  \longrightarrow\mathbf{Pic}^{\mathrm{loc}}(v,C(V,M))
  \longrightarrow\mathbf{Pic}(V)/\mathbb Z[M]
  \longrightarrow0.
\end{equation}
The left term is a vector group, not merely a tangent
space. The sequence is functorial for the hyperplane
restriction under consideration. These are
\citet[Example 4.2, equations (4.2.3)--(4.2.4),
p.~470]{D_Kollar2016LocalPicard}.
Its elementary mechanism is the proper cone contraction
from $\operatorname{Tot}(M^{-1})$:
on the formal line bundle the characteristic-zero
logarithm identifies $1+I$ with $I$, and
$I=\prod_{a>0}M^a$. Restriction to the zero section has
the splitting supplied by the projection to $V$.
Formal existence for this proper contraction gives its line
bundles, and removing the exceptional divisor divides
out by its class $[M]$. In particular this splitting
is compatible with the inclusion $H\hookrightarrow W$ and the
restriction maps. The local Picard-scheme calculation
in the cited result makes these actual vector groups.
For the full section ring $B_H$, this contraction is
the blow-up of the ideal generated by the degree-one
sections; we do not assert that it is the blow-up of
the maximal ideal. Indeed $B_H$ need not be generated
in degree one. In the present example the section-ring
defect is concentrated in degree two and is
$H^1(W,L)\simeq\mathbb C^{m+1}$. On a chart where a
degree-one section $t$ is chosen, the contraction has
ring $\mathcal O(H_t)[t]$, proving this description
directly. Its exceptional divisor and formal normal
bundle are still exactly $(H,L_H^{-1})$.

Also, over an algebraically closed ground field,
\begin{equation}\label{eq:C3-henselian-Picard}
 \mathbf{Pic}^{\mathrm{loc}}(v,C(V,M))(\mathbb C)
   =\operatorname{Pic}
       \bigl(\operatorname{Spec}\mathcal O^{\mathrm{sh}}_{C(V,M),v}
                                                     \setminus\{v\}\bigr).
\end{equation}
This is \citet[Section 3.5, p.~469]{D_Kollar2016LocalPicard}.
It is the point that passes from a local Picard scheme
to the requested strict henselian puncture. We are
not inferring a line-bundle counterexample from a
tangent vector, nor replacing strict henselization
by completion without a comparison theorem.

The cohomology and its restriction can now be calculated
explicitly. Kunneth and Serre duality give
\[
 \begin{aligned}
 H^1(W,\mathcal O_W)&=H^1(C,\mathcal O_C)
                           =H^0(C,K_C)^*,\\
 H^1(W,L)&=H^0(\mathbb P^m,\mathcal O(1))
                            \otimes H^1(C,K_C)
                              \simeq\mathbb C^{m+1},\\
 H^1(W,L^a)&=0\quad(a\ge2),\qquad
 H^2(W,L^a)=0\quad(a\ge0).
 \end{aligned}
\]
Since $z_0,z_1,z_2$ form a basis of $H^0(C,K_C)$,
multiplication by $f$ on the first of these groups is
\[
 H^0(C,K_C)^*\longrightarrow
                H^0(\mathbb P^m,\mathcal O(1)),\qquad
 \lambda\longmapsto\sum_{j=0}^{2}\lambda(z_j)x_j.
\]
It is injective with image
$V_0=\langle x_0,x_1,x_2\rangle\simeq\mathbb C^3$.
The hyperplane exact sequence therefore yields
\[
 H^1(H,L_H)=H^1(W,L)/V_0,
 \qquad H^1(H,L_H^a)=0\quad(a\ge2).
\]
Thus restriction on the vector groups in
\eqref{eq:C3-Kollar-cone} is the quotient
$\mathbb C^{m+1}\to\mathbb C^{m-2}$ with kernel $V_0$.

Finally the projective-bundle formula gives
\[
 \operatorname{Pic}(W)=\mathbb Z[\mathcal O_{\mathbb P^m}(1)]
                              \oplus\operatorname{Pic}(C),
 \qquad
 \operatorname{Pic}(H)=\mathbb Z[\mathcal O_{\mathbb P(E)}(1)]
                              \oplus\operatorname{Pic}(C).
\]
Restriction is an isomorphism between these groups and
carries $L$ to $L_H$; it is the same isomorphism on the
corresponding Picard schemes. Applying
\eqref{eq:C3-Kollar-cone} to $W$ and $H$, with the
compatible projection splittings, shows that the full
kernel on local Picard schemes is precisely the
vector group $V_0$, and that the map is surjective.
Equation \eqref{eq:C3-henselian-Picard} now proves
\eqref{eq:C3-picard-kernel} for actual line bundles.
Every nonzero point of this additive group is
non-torsion. On the regular punctures,
$\operatorname{Pic}=\mathrm{CH}^1$, giving the asserted
refutations.
\end{proof}

\emph{The necessary distinction concerning $J^1$.}
The same calculation refutes injectivity of the
\emph{identity components of the local Picard schemes}:
the displayed vector kernel is already in their
identity components. In this specific example those
components are
\[
 \operatorname{Pic}^0(C)\times\mathbb G_a^{m+1}
       \longrightarrow
 \operatorname{Pic}^0(C)\times\mathbb G_a^{m-2},
\]
with identity on the abelian factor and the displayed
vector-space quotient on the other factor.
However, the algebraic equivalence used in that
local Picard functor permits an etale neighbourhood
of the product with the parameter scheme that varies
with the parameter
\citep[Definition 3.3, clauses (2.a)--(2.c),
pp.~467--468]{D_Kollar2016LocalPicard}.
Ordinary $\mathrm{CH}^1(X)_{\mathrm{alg}}$ uses
families on the fixed scheme $X$. These conventions
must not be identified without proof. Accordingly,
\eqref{eq:C3-picard-kernel} is a refutation of the
full rational-equivalence $A^1$ statements, and of
the local-Picard-identity-component reading of $J^1$;
it is not asserted here to refute
$J^1=\mathrm{CH}^1_{\mathrm{alg}}\otimes\mathbb Q$.
The already visible abelian factor in this example
maps isomorphically.

The missing depth condition is substantive.
Indeed $H^1(W,\mathcal O_W)\ne0$ makes
$\operatorname{depth}A=2$, whereas $f$ is a
nonzerodivisor, so $\operatorname{depth}A/(f)=1$.
This lies outside the characteristic-zero injection
theorem of \citet[Theorem 0.1, p.~833]{D_BhattDeJong2014},
which requires dimension at least four and
$\operatorname{depth}A/(f)\ge2$.
Their Examples 1.31 and 1.33--1.34, pp.~843--844,
also explain why a cohomological or analytic example
requires an additional algebraization argument.
The group-scheme calculation and
\eqref{eq:C3-henselian-Picard} supply the requisite
strict-local argument here, with a smooth projective-bundle
hyperplane section explicitly fixed.
The exact-arithmetic certificate
\texttt{code/verify\_local\_picard.py} verifies the
three affine-chart smoothness and section-generation
certificates, and the contraction/quotient matrices for
$4\le N\le20$. The proof of the assertions for every
$N$ is the uniform calculation above.
